\chapter{Background}
\label{ch:background}

\section{Overview of Conversational Recommender Systems}

Conversational Recommender Systems (CRS) are a subset of RS that are interactive in their nature. CRS  provide users with an opportunity to communicate their preferences and utilise this information to update the recommendations in real time. Recommender chatbots are a subset of CRS which use natural language in order to communicate to users and endeavor to recreate an experience of talking to a human advisor.

CRS employ various approaches to their architecture, but they share common components that form their foundation. Key components that are often found in CRS are:

\begin{itemize}
    \item Dialogue Management System, also known as State Tracker or Conversation Flow Management module, whose goal is to process user actions, update the state of the system and agent’s representation of user preferences, as well as select the next action; the actions commonly include request for information, offering a recommendation, explanation or justification of a recommendation, as well as other forms of response to user’s actions
    \item User Modelling System, which recalculates user’s representation based on the available data
    \item Recommendation Engine, which either uses newly obtained information to filter a precalculated list of recommendations \citep{sun2018conversational} or updates a recommender algorithm which can be a mixture of collaborative filtering, content-based filtering and knowledge based recommender algorithm \citep{sun2018conversational, christakopoulou2016towards} and ranks items in the order of estimated relevance.
    \item Input and Output Processing Modules, which include natural language understanding and generation components
    \item Knowledge Elements, which include item databases and other domain knowledge.
\end{itemize}

Recent survey paper that review the state of the field of CRS acknowledges that there is a wide variation in CRS designs both algorithmically and in terms of user interface \citep{jannach2021survey, pramod2022conversational}. Different conversational recommender system architectures can be compared and classified in a number of ways. The most prominent dimensions are (1) text/speech based vs visual interface based, (2) collaborative filtering vs attribute based and (3) preference-based vs critique based. These distinctions will be discussed in more detail in the following paragraphs.

\textit{User interface}. Some interactive recommender systems personalise web pages \citep{mahmood2014dynamic} and menus in a form of buttons, cards, carousels \citep{ayundhita2019ontology, dalton2018vote} based on customer feedback and behaviour. Others converse with customers purely in natural language form (text or voice) \citep{dalton2018vote} . While having chat bot capability imitates human beings conversing and clearly better approximates real world settings, arguably, pure natural language is not always appropriate, as some product dimensions are easier to visualise than express verbally. In many circumstances, it could be helpful for users to perceive the visual aesthetics of the product on offer or access embedded hyperlinks that link to the pages with additional product information. Many systems successfully combine both natural language and non-verbal interaction capabilities \citep{dalton2018vote}.

\textit{Recommendation system type}. While some conversational recommenders use collaborative filtering approach both for the recommendation modelling and for question formulation \citep{dalton2018vote, zhao2013interactive}, others utilise item attribute information contained in user requests \citep{mahmood2014dynamic, warnestaal2005modeling} or actively elicit such information from customers by asking facet-related questions that mimic filters on a typical e-commerce site \citep{sun2018conversational}. While collaborative information is often used in natural human conversations in domains like movies and so this approach translates well to a conversational recommender system \citep{warnestaal2005modeling}, arguably, in many domains such approach is impossible. This particularly concerns areas where purchases are infrequent (travel, white goods) and individual product titles or models are not descriptive. For instance, it would be unnatural to ask a customer: ‘Name some of the hotels that you liked’. In contrast, natural human interactions between a travel agent and a customer are likely to be dominated by attribute based questions, such as: ‘Where do you want to go?’ ‘Would you like a five star hotel?’ ‘Do you require a hotel with a creche?’ etc. Some systems employ a third approach, where product attributes are mapped onto user functional requirements to mitigate possible user unfamiliarity with technical specifications \citep{ayundhita2019ontology}.

\textit{Information elicitation approach}. Finally, conversational recommender systems employ different strategies for information elicitation. While some systems are designed to formulate questions related to customer preferences or previous purchase history \citep{sun2018conversational}, others apply a trial-and-error strategy, where an algorithm takes several attempts at returning optimal recommendations and asks a user to provide feedback on these recommendations \citep{li2018towards, mcginty2006adaptive}. The latter approach is known as critique recommender system.


\subsection{Conversational Recommender System Training Approaches}

Some conversational systems use rule-based approaches of varying complexity to guide the dialogue flow \citep{ayundhita2019ontology, dalton2018vote, chandrashekara2018fuzzy}. This can be represented as dialogue flow graph \citep{warnestaal2005modeling} or a simple set of rules \citep{dalton2018vote}. A small group of natural language interfaced systems rely on supervised learning only \citep{li2018towards}. Majority of CRS, however, utilise on reinforcement learning algorithms that allow to learn optimal actions based on user behaviour \citep{sun2018conversational, christakopoulou2016towards}. 

In contrast to traditional recommender systems, CRS are interactive in their nature. CRS provide users with an opportunity to communicate their preferences and utilise this information to update user representations in real time. Despite the differences in design and mode of interaction, the main goal of all such systems is to elicit information from users in order to optimise recommendations in a way that increases their relevance. This makes CRS a perfect use case for applying a range of explore-exploit techniques for a number of reasons. 

Firstly, CRS imply continuous interaction with the external environment (a user), which ensures that rewards can be easily operationalised and fed back to the CRS algorithm in real time. This frees reinforcement learning driven CRS from rigid assumptions that characterise systems driven by handcrafted rules \citep{mahmood2009learning} or trained using supervised learning only \citep{li2018towards}, such as that a certain dialogue flow is optimal for all users, or that recommendations should mimic behaviour that users exhibit naturally, an idea exploited by various collaborative filtering and predictive algorithms. Instead, the system behaviour can be tailored to directly optimise performance metrics meaningful in real world settings, such as dialogue success rate \citep{vendrov2020gradient, lei2020estimation, li2021seamlessly}, including business performance indicators like conversion and user satisfaction \citep{mahmood2009improving, mahmood2009learning}.

For example, \cite{mahmood2009improving} show that an adaptive reinforcement learning driven system outperforms handcrafted rules in a real travel planning application by resulting in shorter and quicker system-user interactions, as well as higher conversion and cross-sell rates. Similarly, multiple works demonstrate that fine-tuning pretrained algorithms with reinforcement learning significantly improves the quality of resulting recommendations. This has been shown for supervised learning algorithms \citep{henderson2008hybrid}, collaborative filtering algorithms \citep{sun2018conversational, li2021seamlessly, lei2020estimation} and entropy based methods for preference elicitation \citep{sun2018conversational}. For instance, \cite{sun2018conversational} argue that reinforcement learning driven models are more adaptive and, therefore, are capable of learning at which point it becomes beneficial to stop preference elicitation and display recommendations to the user. Furthermore, \cite{li2021seamlessly} demonstrate that powerful exploration techniques such as Thompson sampling allow to improve the dialogue success rate on a state-of-the-art CRS \citep{lei2020estimation} by 70\%-150\%, with 16\% to 76\% of dialogues ending in a successful recommendation, depending on the dataset. They argue that effective exploration allows to quickly move from a prior if it does not accurately reflect customer needs.

Secondly, reinforcement learning allows CRS to overcome such critical problems as cold start or user taste drift \citep{vendrov2020gradient} by leveraging active exploration techniques. Such systems converge to an optimal set of recommendations over the course of the conversation with a user. For instance, when a user provides negative feedback regarding an item or an attribute, it impacts the rating of all related items and attributes, while further exploration of the same space is discouraged \citep{li2021seamlessly}. Even in complete absence of any a priori information, reinforcement learning systems can efficiently elicit user preferences and produce relevant recommendations in just a few conversational turns. For example, it has been shown that a randomly initialised system in certain circumstances can catch up with a system initialised based on data gathered offline in 13 conversational turns \citep{christakopoulou2016towards}. 

Thirdly, reinforcement learning techniques can be totally data driven, which makes them easily transferable to new application scenarios. In contrast, it has been shown that non-adaptive systems can show unreliable and wildly different performance on new datasets \citep{li2021seamlessly}.


\section{Recommender System Architectures in Conversational Recommender Systems}

Recommender Systems (RS) are an established and commercially acclaimed family of algorithms designed to help users by narrowing down the range of potential options and surfacing the most relevant choices. It has been previously shown that abundance of choice, particularly when choice complexity is high (e.g., number of distinguishing attributes is large), may lead to customer frustration, which, in turn, results in lower conversion in e-commerce settings \citep{greifeneder2010less}. Recommender systems are increasingly adopted by online businesses to mitigate this effect and increase website traffic and margin \citep{zhao2013interactive}. Furthermore, recommenders are successfully applied in non-for-profit contexts in areas such as science and education \citep{chandrashekara2018fuzzy} and urban planning \citep{kucherbaev2017chatbots}.

\subsection{Neural Recommenders}

Neural recommender systems encompass a broad range of algorithms that often extend traditional approaches. Recent survey paper by \cite{sheng2022deep} reviews deep learning based CRS and concludes that while many authors choose traditional matrix factorisation based algorithms, neural based recommenders become more ubiquitous with architectures including autoencoder based models and graph based CNN models.

According to a recent survey by \cite{wu2022survey}, neural recommender systems can be categorized into three main types: collaborative filtering, content-enriched recommendation, and temporal/sequential recommendation. Collaborative filtering leverages user-item interaction data as the primary source of information. Content-enriched recommendation incorporates additional side information, such as user profiles and item knowledge graphs. Temporal/sequential recommendation takes into account contextual information associated with user interactions, including time, location, and past interactions. Each approach offers unique advantages for improving recommendation accuracy and addressing the evolving nature of user preferences.

Collaborative filtering based approaches include: 

\begin{itemize}
    \item extensions of Matrix Factorisation models, where neural network is used to learn weights that represent importance of individual items in predicting the overall user profile;
    \item Autoencoder-based models, which utilize the concept of reconstructing input to improve representation learning;
    \item graph-based models, which allow to mine higher-order connectivity from the user-item graph structure.
\end{itemize}

Content-rich recommendations allow to improve recommendation quality by incorporating additional signal. This group includes the following approaches:

\begin{itemize}
    \item neural extensions of Factorisation Machines, which encode categorical attributes along with users and items;
    \item  neural Natural Language Processing models enable multi-level automatic representation learning of textual content, such as item descriptions or user-item pair connections (e.g., product reviews);
    \item models that embed multimedia content, especially visual content sich as images and videos;
    \item models that embed social network information.
\end{itemize}

Temporal/sequential models address the dynamic nature of users' preferences, as these preferences evolve over time. Instead of solely modeling users' static preferences, temporal/sequential recommendation focuses on capturing users' dynamic preferences and sequential patterns over time.

Neural recommender systems offer remarkable flexibility, allowing for the adaptation of architectures to suit the specific requirements of a system. In particular, sequence models like LSTM-based architectures and attention models are highly suitable for dialogue-based scenarios, where user profile information is accumulated over multiple conversational turns. These models excel at capturing the intricate dynamics of user interactions, enabling more nuanced and effective recommendations in dialogue-driven recommendation systems.
